%%%%%%%%%%%%%%%%%%%%%%%%%%%%%%%%%%%%%%%%%%%%%%%%%%%%%%%%%%%%%%%%%%%%%%
% How to use writeLaTeX: 
%
% You edit the source code here on the left, and the preview on the
% right shows you the result within a few seconds.
%
% Bookmark this page and share the URL with your co-authors. They can
% edit at the same time!
%
% You can upload figures, bibliographies, custom classes and
% styles using the files menu.
%
%%%%%%%%%%%%%%%%%%%%%%%%%%%%%%%%%%%%%%%%%%%%%%%%%%%%%%%%%%%%%%%%%%%%%%

\documentclass[12pt]{article}

\usepackage{sbc-template}

\usepackage{graphicx,url}

%\usepackage[brazil]{babel}   
\usepackage[utf8]{inputenc}  

     
\sloppy

\title{Exercício 4 do Grau B Sistemas Distribuidos}

\author{João Roberto Lemos Fidellis\inst{1}, Christian Aguiar Plentz\inst{2}, 
Julian Rafael de Souza\inst{3}}


\address{Ciência da Computação - Universidade do Vale do Rio dos Sinos (UNISINOS)\\
  Caixa Postal 275 -- 93.022-750 -- São Leopoldo -- RS -- Brazil
\email{\{plentzchristian, jrfidellis, julian93\}@edu.unisinos.br}
}

\begin{document} 

\maketitle

\section{Apache CXF para Web Services}

O Apache CXF é um framework escrito em Java para criação e consumo de Web Services baseado em SOAP e REST. Possui a abordagem Code-First para implementação o Desenvolvedor precisa anotar suas Funções desejadas com a anotação \textit{@WebService} e \textit{@WebMethod} vide o exemplo abaixo.

\begin{verbatim}
import jakarta.jws.WebService;
import jakarta.jws.WebMethod;

@WebService
public class CalculadoraService {
	@WebMethod
	public int somar(int a, int b){
		return a + b;
	}
}
\end{verbatim}

Para criar o wsdl dessa classe basta usar o  \\ 
\texttt{Endpoint.publish(ENDERECO\_E\_PORTA, New CalculadoraService())} 
durante seu fluxo de execucao para criar um servidor de wsdl, 
se quiser pode se testar o servidor para ver se está dando retorno com \\
\texttt{curl -v "http://localhost:8080/calculadora?wsdl"}

Para poder gerar o código para implementar o cliente, se pode gerar o código automaticamente
com o programa wsdl2java vide o exemplo abaixo. \\
\texttt{wsdl2java -d . -p client.wsdl \$ENDERECO\_E\_PORTA }

Com isso vai se gerar vários arquivos java relacionados a classe
CalculadoraService para o desenvolvedor poder implementar.

\begin{verbatim}
 $ ls client/wsdl 
CalculadoraService.java         ObjectFactory.java  Somar.java
CalculadoraServiceService.java  package-info.java   SomarResponse.java
\end{verbatim}

Exemplo de uso dentro do cliente:

\begin{verbatim}
import client.wsdl.CalculadoraServiceService;
import client.wsdl.CalculadoraService;

public class Client {
	public static void main(String[] args) {

		try {
			var service = new CalculadoraServiceService();
			var port = service.getCalculadoraServicePort();

			int resultado = port.somar(15, 20);

			System.out.printf("resultado da chamada de somar(15, 20) no Servidor %d", resultado);
		} catch (Exception e) {
			e.printStackTrace();
		}
	}
}

\end{verbatim}

\end{document}

